\chapter{METODOLOGI DAN RANCANGAN PENELITIAN}

\section{Metodologi Penelitian}

Penelitian ini menggunakan pendekatan pengembangan prototipe (\textit{prototyping}) secara iteratif. Pendekatan ini sesuai untuk membangun sistem yang komponennya perlu diuji bersama, yaitu agen pembelajaran, pengirim permintaan HTTP, penganalisis respons, dan mekanisme verifikasi temuan. Setiap iterasi mencakup perancangan, implementasi, pengujian pada lingkungan lokal, serta perbaikan berdasarkan hasil pengamatan \cite{pressman2014}.

Tahapan penelitian diringkas pada Gambar \ref{fig:alur_metodologi_penelitian}. Sistem diuji menggunakan dua agen, yaitu Q-Learning tabular dan Deep Q-Network (DQN). Keduanya berinteraksi dengan lingkungan dan ruang aksi yang sama, tetapi memakai representasi nilai serta jadwal eksplorasi yang sesuai dengan algoritmanya.

\begin{figure}[H]
\centering
\begin{tikzpicture}[
    node distance=0.65cm,
    process/.style={
        draw, rectangle, rounded corners,
        minimum width=11cm, minimum height=0.8cm,
        align=center, fill=blue!10, font=\small\sffamily
    },
    arrow/.style={-{Stealth}, thick}
]
    \node[process] (n1) {Studi literatur dan identifikasi kebutuhan penelitian};
    \node[process, below=of n1] (n2) {Penetapan target, endpoint, payload, state, reward, dan anggaran eksperimen};
    \node[process, below=of n2] (n3) {Implementasi klien HTTP, analyzer, verifikator, agen Q-Learning, dan agen DQN};
    \node[process, below=of n3] (n4) {Penyelidikan baseline per endpoint dan pelatihan masing-masing agen};
    \node[process, below=of n4] (n5) {Benchmark detektor dan pemindaian terkontrol OWASP ZAP};
    \node[process, below=of n5] (n6) {Pengolahan log, perhitungan metrik, dan pembahasan hasil};
    \draw[arrow] (n1) -- (n2);
    \draw[arrow] (n2) -- (n3);
    \draw[arrow] (n3) -- (n4);
    \draw[arrow] (n4) -- (n5);
    \draw[arrow] (n5) -- (n6);
\end{tikzpicture}
\caption{Tahapan Metodologi Penelitian}
\label{fig:alur_metodologi_penelitian}
\end{figure}

\section{Alat dan Bahan yang Digunakan}

\subsection{Perangkat Lunak}

Perangkat lunak yang digunakan dalam penelitian dirangkum pada Tabel \ref{tab:perangkat_lunak}.

\begin{table}[H]
\centering
\caption{Perangkat Lunak yang Digunakan}
\label{tab:perangkat_lunak}
\begin{tabular}{|c|p{3.3cm}|p{8.4cm}|}
\hline
\textbf{No.} & \textbf{Komponen} & \textbf{Kegunaan} \\ \hline
1 & Python & Menjalankan agen, klien HTTP, analisis respons, verifikasi, dan pengolahan hasil. \\ \hline
2 & NumPy & Mendukung operasi numerik pada implementasi agen dan pengolahan state. \\ \hline
3 & PyTorch & Menjalankan jaringan saraf, replay buffer, dan pembaruan target pada DQN. \\ \hline
4 & Requests & Mengirim permintaan HTTP ke endpoint OWASP Juice Shop. \\ \hline
5 & Playwright Chromium & Memeriksa apakah payload XSS benar-benar dieksekusi pada DOM halaman frontend. \\ \hline
6 & OWASP Juice Shop v20.2.0 & Aplikasi web lokal yang menjadi lingkungan target pengujian \cite{owaspJuiceShop}. \\ \hline
7 & OWASP ZAP & Pemindai pembanding yang dikendalikan melalui Core API untuk cakupan dua endpoint \cite{owaspZapApi}. \\ \hline
\end{tabular}
\end{table}

\subsection{Perangkat Keras dan Lingkungan}

Juice Shop dijalankan pada komputer lokal dan diakses melalui \texttt{http://localhost:3000}. Agen dan layanan target berada pada lingkungan yang sama agar lalu lintas eksperimen tidak bergantung pada jaringan eksternal. Browser Chromium digunakan oleh \texttt{BrowserVerifier} untuk pengujian eksekusi DOM. Spesifikasi CPU, RAM, sistem operasi, dan akselerator komputasi yang benar-benar digunakan dicatat dari perangkat eksperimen.

\section{Arsitektur Sistem}

Arsitektur sistem terdiri atas agen RL, pemetaan aksi ke payload dan endpoint, klien HTTP, aplikasi target, analyzer respons, verifikator, serta komponen reward dan pencatatan. Agen menerima state, memilih indeks aksi, lalu mengirimkannya ke pembangkit payload. Klien HTTP meneruskan payload ke Juice Shop. Respons dianalisis untuk membentuk kandidat, tetapi kandidat hanya menjadi temuan setelah melewati gerbang verifikasi. State berikutnya dan reward kemudian digunakan untuk pembaruan agen, sedangkan seluruh transaksi dicatat untuk evaluasi.

\begin{figure}[H]
\centering
\begin{tikzpicture}[
    node distance=0.55cm,
    module/.style={
        draw, rectangle, rounded corners,
        minimum width=10.5cm, minimum height=0.75cm,
        align=center, fill=blue!7, font=\small
    },
    arrow/.style={-{Stealth}, thick}
]
    \node[module] (agent) {Agen RL: Q-Learning atau DQN};
    \node[module, below=of agent] (mapping) {Pemetaan aksi: indeks endpoint dan indeks payload};
    \node[module, below=of mapping] (http) {Klien HTTP: penyusunan dan pengiriman permintaan};
    \node[module, below=of http] (target) {OWASP Juice Shop lokal};
    \node[module, below=of target] (analyzer) {Response Analyzer: fitur respons, state, dan kandidat temuan};
    \node[module, below=of analyzer] (verifier) {Gerbang verifikasi: kontrol diferensial atau eksekusi DOM};
    \node[module, below=of verifier] (update) {Reward, pembaruan agen, dan pencatatan log};
    \draw[arrow] (agent) -- (mapping);
    \draw[arrow] (mapping) -- (http);
    \draw[arrow] (http) -- (target);
    \draw[arrow] (target) -- (analyzer);
    \draw[arrow] (analyzer) -- (verifier);
    \draw[arrow] (verifier) -- (update);
    \draw[arrow] (update.west) -- ++(-1.3,0) |- node[pos=0.25,left,align=center,font=\scriptsize] {state\\dan reward} (agent.west);
\end{tikzpicture}
\caption{Arsitektur Sistem Pengujian Berbasis Reinforcement Learning}
\label{fig:arsitektur_sistem_rl}
\end{figure}

\section{Rancangan Lingkungan dan Ruang Aksi}

\subsection{Lingkungan Pengujian}

Lingkungan pengujian adalah OWASP Juice Shop v20.2.0 yang berjalan secara lokal. Sistem menggunakan delapan endpoint tetap yang telah diperiksa dapat diakses. Daftar endpoint, metode HTTP, parameter, serta permukaan pengujiannya disajikan pada Tabel \ref{tab:endpoint_target}.

\begin{table}[H]
\centering
\caption{Endpoint yang Digunakan dalam Eksperimen}
\label{tab:endpoint_target}
\begin{tabular}{|c|p{5.2cm}|c|p{3.2cm}|p{2.3cm}|}
\hline
\textbf{Indeks} & \textbf{Endpoint} & \textbf{Metode} & \textbf{Input} & \textbf{Surface} \\ \hline
0 & \texttt{/rest/products/search?q=} & GET & Parameter \texttt{q} & SQL \\ \hline
1 & \texttt{/ftp/} & GET & Parameter/path berkas & File \\ \hline
2 & \texttt{/rest/products/1/reviews} & GET & Parameter \texttt{id} & Plain \\ \hline
3 & \texttt{/rest/track-order/} & GET & Parameter/path pesanan & Plain \\ \hline
4 & \texttt{/rest/user/login} & POST & JSON & SQL \\ \hline
5 & \texttt{/rest/products/1/reviews} & PUT & JSON & Plain \\ \hline
6 & \texttt{/api/Feedbacks} & POST & JSON dan CAPTCHA & Plain \\ \hline
7 & \texttt{/api/Challenges} & GET & Parameter \texttt{q} & Plain \\ \hline
\end{tabular}
\end{table}

Nilai \textit{surface} menyatakan jenis pemrosesan input yang relevan pada endpoint, bukan label bahwa endpoint tersebut pasti rentan. Surface SQL hanya dipakai pada endpoint indeks 0 dan 4, sedangkan surface file dipakai pada indeks 1. Endpoint lainnya diperlakukan sebagai surface plain. Pembatasan ini mencegah respons yang bentuknya tampak serupa diklasifikasikan sebagai kerentanan yang tidak sesuai dengan cara kerja endpoint.

Sistem memuat 15 payload yang telah disusun ke dalam tiga kelompok: enam payload SQL injection, empat payload XSS, dan lima payload fuzzing. Payload serta pemetaan surface merupakan prior yang disusun peneliti dengan pengetahuan tentang struktur Juice Shop. Hal ini perlu diungkapkan sebagai \textit{human prior}; agen tetap harus memilih aksi yang tepat, dan indikasi awal tidak dihitung sebagai temuan sebelum melewati verifikasi.

Untuk verifikasi XSS, endpoint indeks 0 dipetakan ke halaman \texttt{/\#/search?q=} dan endpoint indeks 3 ke halaman \texttt{/\#/track-result?id=}. Pemetaan ini digunakan oleh verifikator browser dan tidak mengubah ruang aksi HTTP.

\subsection{Ruang Aksi}

Satu aksi adalah pasangan endpoint dan payload. Dengan delapan endpoint dan 15 payload, ruang aksi berisi 120 kombinasi. Indeks aksi \(a\) dikodekan sebagai:

\begin{equation}
i_{\text{endpoint}} = \left\lfloor \frac{a}{15} \right\rfloor,
\qquad
j_{\text{payload}} = a \bmod 15,
\qquad
a \in \{0,1,\ldots,119\}.
\label{eq:action_encoding}
\end{equation}

Dengan demikian, agen memilih satu kombinasi endpoint-payload pada setiap episode. Untuk permintaan POST dan PUT, payload dikirimkan melalui badan JSON sesuai kebutuhan endpoint. Endpoint umpan balik indeks 6 memerlukan CAPTCHA; klien mengambil nilai CAPTCHA melalui endpoint terkait dan melengkapi nilai bawaan yang dibutuhkan oleh permintaan.

\section{Rancangan Agen Reinforcement Learning}

\subsection{State}

State merangkum informasi respons yang digunakan agen untuk menentukan aksi berikutnya. Implementasi menggunakan tuple delapan komponen seperti pada Tabel \ref{tab:state_components}.

\begin{table}[H]
\centering
\caption{Komponen State yang Digunakan}
\label{tab:state_components}
\begin{tabular}{|c|p{3.2cm}|p{3.6cm}|p{5.3cm}|}
\hline
\textbf{No.} & \textbf{Komponen} & \textbf{Representasi} & \textbf{Kategorisasi} \\ \hline
1 & Status HTTP & Kategori status & \texttt{2xx}, \texttt{3xx}, \texttt{4xx}, \texttt{5xx}, atau \texttt{error}. \\ \hline
2 & Panjang respons & Kategori 0--3 & Panjang teks respons terdekode: \(<1500\), \(1500\)--\(<2500\), \(2500\)--\(<3500\), atau \(\geq 3500\) karakter. \\ \hline
3 & Indikator error & Biner & Kecocokan pola pesan error basis data. \\ \hline
4 & Indikator refleksi & Biner & Payload atau token bermakna terdeteksi pada respons. \\ \hline
5 & Indikator kerentanan & Biner & Bernilai 1 jika temuan tetap terklasifikasi setelah gerbang verifikasi. \\ \hline
6 & Indeks endpoint & Bilangan bulat 0--7 & Menunjukkan endpoint yang menjadi sasaran aksi. \\ \hline
7 & Kategori waktu respons & Kategori 0--2 & Cepat \(<0{,}5\) detik, sedang \(0{,}5\)--\(<2\) detik, atau lambat \(\geq 2\) detik. \\ \hline
8 & Kategori tipe konten & Kategori 0--4 & JSON, HTML, teks, XML, atau tipe lainnya. \\ \hline
\end{tabular}
\end{table}

Panjang respons pada implementasi dihitung dari teks respons yang telah didekode, bukan dari ukuran byte mentah. Kategori panjang menggunakan batas 1.500, 2.500, dan 3.500 karakter. Baseline perubahan respons dihitung per endpoint berdasarkan probe netral. Jika baseline endpoint tidak tersedia, nilai fallback yang digunakan adalah 2.200 karakter. Perubahan panjang dianggap signifikan apabila deviasi relatif terhadap baseline mencapai 15\%.

State awal yang digunakan kedua agen adalah:

\begin{equation}
s_0 = (\texttt{2xx}, 1, 0, 0, 0, 0, 0, 0).
\end{equation}

State setelah suatu aksi menjadi masukan pada iterasi berikutnya. Pada tiap nomor episode, agen melakukan satu aksi payload dan menghasilkan satu transisi pembelajaran.

\subsection{Action Masking}

Action masking membatasi pemilihan ulang terhadap aksi yang telah menghasilkan kerentanan terverifikasi. Aksi hanya dimasukkan ke dalam himpunan \texttt{discovered\_actions} setelah lolos kontrol-diferensial atau verifikasi DOM. Aksi yang gagal, tidak menghasilkan temuan, atau ditolak oleh verifikator tidak dimask secara permanen. Dengan kebijakan ini, agen tetap dapat memilih aksi yang sebelumnya tidak berhasil dan hasil discovery tidak ditentukan oleh penghapusan sistematis seluruh aksi yang belum berhasil.

\subsection{Reward}

Reward memberi umpan balik atas hasil setiap aksi. Reward dasar dan bonus kategori panjang dirangkum pada Tabel \ref{tab:reward_scheme}.

\begin{table}[H]
\centering
\caption{Skema Reward dan Bonus Panjang Respons}
\label{tab:reward_scheme}
\begin{tabular}{|p{5.1cm}|c|p{6.6cm}|}
\hline
\textbf{Kondisi} & \textbf{Reward dasar} & \textbf{Keterangan} \\ \hline
Kerentanan terverifikasi & \(+50\) & Diberikan setelah bukti eksploitasi lolos gerbang verifikasi. \\ \hline
Anomali atau indikasi yang belum terverifikasi & \(+10\) & Mencakup kandidat yang ditolak atau respons yang menunjukkan error, refleksi, maupun perubahan signifikan. \\ \hline
Respons normal & \(-1\) & Permintaan berhasil, tetapi tidak menghasilkan temuan atau anomali yang relevan. \\ \hline
Kesalahan permintaan/koneksi & \(-5\) & Permintaan gagal diselesaikan atau koneksi ke target mengalami error. \\ \hline
\end{tabular}
\end{table}

Bonus panjang respons ditambahkan berdasarkan kategori panjang: kategori 0 memperoleh \(+0\), kategori 1 memperoleh \(+0\), kategori 2 memperoleh \(+5\), dan kategori 3 memperoleh \(+10\). Reward akhir dirumuskan sebagai:

\begin{equation}
R_t = R_{\text{dasar},t} + B_{\text{panjang},t}.
\label{eq:reward}
\end{equation}

Reward \(+50\) hanya berhubungan dengan temuan yang lolos verifikasi. Indikasi awal seperti refleksi di badan respons HTTP tidak dengan sendirinya membuktikan bahwa payload berhasil dieksploitasi.

\subsection{Q-Learning Tabular}

Q-Learning memperbarui estimasi nilai untuk setiap pasangan state dan aksi. Pembaruan menggunakan persamaan berikut \cite{watkins1992qlearning,sutton2018}:

\begin{equation}
Q(s_t,a_t) \leftarrow Q(s_t,a_t)
+ \alpha \left[
R_t + \gamma \max_{a'} Q(s_{t+1},a') - Q(s_t,a_t)
\right],
\label{eq:qlearning}
\end{equation}

dengan \(s_t\) sebagai state saat ini, \(a_t\) sebagai aksi yang dipilih, \(R_t\) sebagai reward, \(s_{t+1}\) sebagai state berikutnya, \(\alpha\) sebagai laju pembelajaran, dan \(\gamma\) sebagai faktor diskonto. Q-Table diinisialisasi dengan nilai optimistis \(5{,}0\) untuk setiap aksi yang belum memiliki estimasi.

\subsection{Deep Q-Network}

DQN mengaproksimasi nilai \(Q(s,a)\) menggunakan jaringan saraf tiruan dengan dua lapisan tersembunyi berukuran 64 neuron dan aktivasi ReLU. State delapan komponen diubah menjadi vektor 28 dimensi: status HTTP (5), kategori panjang (4), tiga indikator biner, endpoint (8), waktu respons (3), dan tipe konten (5). Lapisan keluaran menghasilkan satu nilai Q untuk masing-masing dari 120 aksi. Rancangan DQN mengikuti penggunaan replay buffer dan jaringan target untuk membantu menstabilkan pembelajaran \cite{mnih2015dqn}.

Target pembelajaran DQN dihitung sebagai:

\begin{equation}
y_t = R_t + \gamma(1-d_t)
\max_{a'} Q_{\theta^-}(s_{t+1},a'),
\label{eq:dqn_target}
\end{equation}

dengan \(d_t\) sebagai penanda transisi terminal dan \(Q_{\theta^-}\) sebagai jaringan target. Jaringan utama diperbarui menggunakan galat kuadrat antara \(y_t\) dan \(Q_\theta(s_t,a_t)\). Pengalaman transisi disimpan dalam replay buffer dan diambil dalam mini-batch.

\subsection{Hyperparameter}

Kedua agen menggunakan anggaran 1.585 episode. Setiap episode memuat satu aksi payload. Q-Learning dan DQN menggunakan jadwal eksplorasi yang berbeda karena keduanya memiliki karakteristik pembelajaran yang berbeda. Nilai yang digunakan diringkas dalam Tabel \ref{tab:hyperparameter}.

\begin{table}[H]
\centering
\caption{Hyperparameter Eksperimen Q-Learning dan DQN}
\label{tab:hyperparameter}
\begin{tabular}{|p{4cm}|c|c|p{5cm}|}
\hline
\textbf{Hyperparameter} & \textbf{Q-Learning} & \textbf{DQN} & \textbf{Keterangan} \\ \hline
Learning rate & \(0{,}1\) & \(0{,}001\) & Laju pembaruan Q-Table atau bobot jaringan. \\ \hline
Discount factor \(\gamma\) & \(0{,}9\) & \(0{,}9\) & Bobot estimasi reward masa depan. \\ \hline
Epsilon awal & \(1{,}0\) & \(1{,}0\) & Probabilitas awal memilih aksi eksplorasi. \\ \hline
Epsilon minimum & \(0{,}1\) & \(0{,}1\) & Batas bawah probabilitas eksplorasi. \\ \hline
Epsilon decay & \(0{,}90\) & \(0{,}9965\) & Faktor peluruhan per episode. \\ \hline
Inisialisasi Q optimistis & \(5{,}0\) & -- & Nilai awal aksi yang belum memiliki estimasi pada Q-Table. \\ \hline
Ukuran lapisan tersembunyi & -- & \(64\) & Banyak neuron pada setiap lapisan tersembunyi DQN. \\ \hline
Ukuran mini-batch & -- & \(32\) & Banyak transisi dalam satu pembaruan DQN. \\ \hline
Kapasitas replay buffer & -- & \(10.000\) & Jumlah maksimum transisi yang disimpan. \\ \hline
Interval pembaruan target & -- & \(10\) langkah & Frekuensi sinkronisasi jaringan target. \\ \hline
Episode per agen & \(1.585\) & \(1.585\) & Anggaran penuh; pelatihan tidak dihentikan lebih awal. \\ \hline
Random seed & \(42\) & \(42\) & Seed yang ditetapkan pada skrip pelatihan. \\ \hline
\end{tabular}
\end{table}

Nilai peluruhan DQN tidak disamakan dengan Q-Learning. Peluruhan \(0{,}90\) akan menurunkan epsilon Q-Learning hingga batas minimum dalam tahap awal pelatihan, sedangkan DQN menggunakan peluruhan \(0{,}9965\) agar eksplorasi berlangsung lebih lama. Karena itu, perbandingan kedua agen dilakukan dengan melaporkan konfigurasi dan hasil aktual masing-masing, bukan dengan menganggap semua hyperparameter harus identik.

\section{Pengumpulan Data dan Verifikasi Temuan}

\subsection{Penyelidikan Baseline per Endpoint}

Sebelum pelatihan, sistem mengirim satu permintaan netral berisi \texttt{hello} ke masing-masing dari delapan endpoint. Panjang respons yang diperoleh disimpan sebagai baseline endpoint tersebut. Analyzer menggunakan baseline yang sesuai saat menghitung perubahan panjang respons; apabila probe baseline gagal, sistem menggunakan fallback 2.200 karakter. Permintaan baseline dilaporkan terpisah dari anggaran 1.585 aksi payload per agen.

\subsection{Deteksi Awal dan Gerbang Verifikasi}

Analyzer menghasilkan kandidat berdasarkan bukti respons dan surface endpoint. SQL injection dapat ditandai berdasarkan pesan error basis data yang relevan atau perubahan respons yang signifikan pada endpoint bersurface SQL. Kandidat file disclosure memerlukan bukti isi berkas atau bukti pemotongan ekstensi pada endpoint file. Kandidat fuzzing memerlukan kondisi respons yang sesuai, termasuk respons 5xx yang menggemakan probe. Refleksi HTTP saja tidak cukup untuk menyatakan XSS.

Kandidat SQL injection dan fuzzing diverifikasi menggunakan kontrol diferensial. Sistem mengirim probe netral \texttt{hello} serta, jika berbeda dari payload asli, varian payload yang dinetralkan. Kandidat diterima apabila kontrol tidak menghasilkan kelas kerentanan yang sama. Jika koneksi gagal atau hasil kontrol mereproduksi kelas yang sama, kandidat tidak dinyatakan terverifikasi. Jumlah permintaan kontrol yang benar-benar dikirim dicatat, sehingga dapat berjumlah satu atau dua untuk satu kandidat.

XSS diverifikasi melalui Playwright Chromium pada halaman frontend yang sesuai. Verifikator menginstrumentasi bentuk tag dan event payload, lalu mengarahkan handler ke canary pada halaman yang dimuat. Kandidat diterima apabila canary terpicu oleh payload uji dan tidak terpicu oleh varian kontrol yang dinetralkan. Payload yang hanya muncul pada respons JSON tidak dianggap sebagai bukti eksekusi. Verifikator melakukan pemanasan satu kali dan menggunakan polling terbatas untuk memeriksa canary agar hasil tidak bergantung pada jeda tetap \cite{playwrightPageApi}.

Suatu temuan dinyatakan terkonfirmasi hanya jika jenis kerentanannya termasuk SQLi, XSS, atau fuzzing dan kolom \texttt{verified} bernilai benar. Kandidat yang ditolak, maupun episode normal, tidak dihitung sebagai temuan terkonfirmasi.

\subsection{Pencatatan Data}

Setiap episode dicatat dalam berkas log terpisah, yaitu \texttt{log\_ql.csv} untuk Q-Learning dan \texttt{log\_dqn.csv} untuk DQN. Ringkasan disimpan masing-masing pada \texttt{summary\_ql.json} dan \texttt{summary\_dqn.json}. Data yang direkam mencakup indeks episode dan aksi, endpoint, payload, status HTTP, panjang dan waktu respons, kategori state, indikator analisis, klasifikasi kerentanan, status verifikasi, reward, serta state sebelum dan sesudah aksi. Ringkasan eksperimen menyimpan jumlah probe baseline, permintaan payload, permintaan kontrol, navigasi browser, temuan terverifikasi, dan metrik lainnya.

Anggaran 1.585 adalah jumlah permintaan payload pelatihan untuk setiap agen, bukan jumlah seluruh transaksi jaringan. Probe baseline dan permintaan kontrol menambah lalu lintas HTTP, sedangkan navigasi browser dilaporkan secara terpisah. Pembedaan ini diperlukan ketika membandingkan jumlah permintaan dengan pemindai lain.

\section{Rancangan Evaluasi}

\subsection{Metrik Pelatihan}

Kinerja agen diringkas dengan metrik berikut:

\begin{enumerate}
    \item \textbf{Temuan terverifikasi}: jumlah episode dengan kelas kerentanan yang lolos verifikasi.
    \item \textbf{Lokasi kerentanan unik}: jumlah pasangan endpoint dan kelas kerentanan yang berbeda. Beberapa payload pada endpoint dan kelas yang sama dihitung sebagai beberapa aksi terverifikasi, tetapi satu lokasi unik.
    \item \textbf{Vuln hit rate}: proporsi episode pelatihan yang menghasilkan temuan terverifikasi:
    \begin{equation}
    \text{Vuln Hit Rate} =
    \frac{N_{\text{episode dengan temuan terverifikasi}}}
    {N_{\text{episode pelatihan}}}.
    \end{equation}
    Metrik ini bukan precision klasifikasi.
    \item \textbf{Requests per vulnerability}: jumlah episode aksi payload dibagi jumlah temuan terverifikasi. Jika tidak ada temuan, nilainya tidak terhingga.
    \item \textbf{Cumulative reward}: jumlah reward yang diterima sepanjang pelatihan, termasuk bonus panjang respons.
    \item \textbf{Waktu pelatihan dan waktu menuju temuan}: durasi eksekusi serta nomor episode temuan pertama dan episode saat temuan terakhir diperoleh.
    \item \textbf{Permintaan tambahan}: jumlah probe baseline dan permintaan verifikasi; navigasi browser dilaporkan terpisah.
\end{enumerate}

Metrik hasil pelatihan bersifat \textit{in-sample}, karena dihitung dari episode yang juga digunakan agen untuk belajar. Eksperimen yang dilaporkan menggunakan satu random seed. Oleh sebab itu, hasil tidak digunakan untuk menyatakan signifikansi statistik atau kemampuan generalisasi ke aplikasi lain.

\subsection{Benchmark Detektor}

Benchmark detektor dijalankan terpisah tanpa agen RL. Analyzer mengevaluasi 120 kombinasi endpoint-payload dan membandingkan klasifikasinya dengan label \textit{ground truth} yang dikurasi manual. Label terdiri atas 14 aksi positif, 60 aksi negatif yang diketahui aman, dan 46 aksi berstatus \textit{unknown}. Aksi \textit{unknown} tidak dimasukkan ke perhitungan positif maupun negatif.

Precision, recall, dan F1 hanya digunakan untuk benchmark detektor:

\begin{equation}
\text{Precision} = \frac{TP}{TP+FP}, \qquad
\text{Recall} = \frac{TP}{TP+FN}, \qquad
F1 = 2 \cdot \frac{\text{Precision}\cdot\text{Recall}}
{\text{Precision}+\text{Recall}}.
\end{equation}

Pelabelan manual dan proses verifikasi menggunakan bukti teknis yang sejenis, yaitu kontrol diferensial untuk sebagian kelas dan eksekusi DOM untuk XSS. Dengan demikian, benchmark ini mengukur konsistensi analyzer terhadap ground truth terkurasi dan false positive pada aksi yang diketahui negatif; hasilnya bukan bukti akurasi universal atau detektor yang sempurna.

\subsection{Pemindaian Pembanding OWASP ZAP}

Pemindaian terkontrol OWASP ZAP mencakup dua endpoint bersama: \texttt{/rest/products/search?q=} dan \texttt{/ftp/}. Skrip pengendali mengakses URL, menjalankan spider, lalu menjalankan active scan dengan opsi \texttt{recurse=false}. Status dan riwayat pesan dibaca melalui Core API; skrip tidak memulai pemindaian pasif sebagai tahap terpisah \cite{owaspZapApi}.

Alert mentah dikelompokkan berdasarkan kategori. Untuk perbandingan temuan aktif, hanya alert SQL injection, XSS, atau path traversal dengan risiko High atau Medium yang dihitung. Alert informasional dan alert konfigurasi tidak diperlakukan sebagai temuan eksploitasi aktif. Versi ZAP dicatat dari API pada saat pemindaian dan harus sama dengan versi yang dilaporkan pada Bab 4.

Hasil Q-Learning juga dihitung ulang untuk endpoint indeks 0 dan 1 agar cakupan endpoint sesuai dengan pemindaian ZAP. Kesamaan cakupan endpoint tidak menyamakan metode: Q-Learning menghitung aksi payload, sedangkan jumlah pesan ZAP mencakup akses URL, spider, dan active scan. Karena satuan permintaan serta strategi pemindaian berbeda, rasio request per temuan disajikan sebagai perbandingan deskriptif pada konfigurasi eksperimen ini, bukan sebagai bukti keunggulan umum salah satu alat.

\subsection{Prosedur Eksperimen}

Urutan eksperimen adalah sebagai berikut:

\begin{enumerate}
    \item Menjalankan OWASP Juice Shop v20.2.0 pada alamat lokal dan memastikan delapan endpoint dapat diakses.
    \item Mengirim probe baseline netral ke setiap endpoint dan menyimpan baseline respons.
    \item Menginisialisasi Q-Learning dan DQN secara terpisah dengan seed 42 serta konfigurasi masing-masing.
    \item Menjalankan 1.585 episode untuk setiap agen. Setiap episode memilih satu aksi dari ruang 120 kombinasi, mengirim payload, menganalisis respons, memverifikasi kandidat, menghitung reward, memperbarui agen, lalu menyimpan log.
    \item Menjalankan benchmark analyzer pada 120 aksi menggunakan label ground truth dan menjalankan pemindaian ZAP terkontrol pada dua endpoint yang telah ditetapkan.
    \item Menghitung metrik dari log, hasil benchmark, dan laporan ZAP, lalu membahas cakupan serta keterbatasan setiap perbandingan pada Bab 4.
\end{enumerate}

Seluruh pengujian diarahkan hanya ke instance Juice Shop lokal yang menjadi target penelitian. Payload, daftar endpoint, dan pengaturan pemindai dibatasi pada lingkungan tersebut.

\section{Jadwal Pelaksanaan}

Jadwal pelaksanaan penelitian dirangkum berdasarkan urutan kegiatan pada Tabel \ref{tab:jadwal_pelaksanaan}. Rentang bulan menunjukkan rencana relatif pelaksanaan kegiatan.

\begin{table}[H]
\centering
\caption{Jadwal Pelaksanaan Penelitian}
\label{tab:jadwal_pelaksanaan}
\begin{tabular}{|p{6.4cm}|c|c|c|c|c|}
\hline
\textbf{Aktivitas} & \textbf{1} & \textbf{2} & \textbf{3} & \textbf{4} & \textbf{5} \\ \hline
Studi literatur & X & X &  &  &  \\ \hline
Analisis kebutuhan dan penetapan ruang lingkup & X &  &  &  &  \\ \hline
Perancangan sistem dan ruang aksi &  & X & X &  &  \\ \hline
Implementasi agen dan modul verifikasi &  &  & X & X &  \\ \hline
Pengujian, benchmark, dan pengumpulan hasil &  &  &  & X & X \\ \hline
Analisis hasil dan penyusunan laporan & X & X & X & X & X \\ \hline
\end{tabular}
\end{table}
